\section{The Harmonic Cochain Problem}

This section states the central search problem and the hidden-operator model in
which it is used.

\begin{definition}[Harmonic Cochain Problem]\label{def:hcp}
  The \emph{Harmonic Cochain Problem} (HCP) asks, given public
  cochain/class-level data associated to a Hodge system, to recover the
  harmonic cochain determined by the operative Laplacian kernel.
\end{definition}

In the \HE{} public-sample model, the visible cochain has the form
\[
  b = h + e,
  \qquad
  h \in \ker(\Delta),
  \qquad
  e \in \image(\delta) \oplus \image(\delta^\ast).
\]
Here \(h\) is the harmonic cochain, equivalently the global-section component,
and \(e\) is Hodge-structured noise. The hidden-operator condition is that the
public data does not determine a unique \(\Delta\) or \(\ker(\Delta)\).

\subsection{Public And Hidden Data}

\begin{definition}[Public observation]\label{def:public-observation}
  A public observation is the data exposed to a verifier or adversary before
  choosing a hidden representative. In the abstract model it consists of a
  public class or shadow value \(p\), together with cochain-facing data \(b\).
  In the \HE{} instantiation, this public layer is represented by public
  shadows, cochain commitments, topology hashes, oracle judgments, and
  proof-carrying artifacts.
\end{definition}

The public observation is intentionally not the hidden Hodge system. It is a
classifier of a public equivalence class of possible hidden representatives.
This is the point at which the fibrational and homotopical semantics enter the
model: public equality is equality after applying a public classifier, not
strict equality of hidden representatives. Section~\ref{sec:mathematical-framework}
develops this semantic layer.

\begin{definition}[Hidden Hodge witness]\label{def:hidden-hodge-witness}
  A hidden Hodge witness for a public observation \((p,b)\) consists of
  hidden operator data \(\Delta\), a harmonic cochain \(h\), and
  Hodge-structured noise \(e\), such that
  \[
    \Delta \in \Hidden(p),
    \qquad
    h \in \ker(\Delta),
    \qquad
    e \in \image(\delta) \oplus \image(\delta^\ast),
    \qquad
    b = h + e.
  \]
  Here \(\Hidden(p)\) denotes the fiber of hidden Hodge--Laplacian
  representatives compatible with the public class value \(p\).
\end{definition}

Concretely, the hidden data includes the fibration representative, restriction
maps, the induced Laplacian, its harmonic kernel, and private witness/cochain
data. The public data may bind to this hidden data through commitments or
proof-carrying artifacts, but it does not present the operative kernel itself.

\begin{definition}[HCP solution]\label{def:hcp-solution}
  A solution to an HCP instance \((p,b)\) is a cochain \(h'\) that is the
  harmonic component selected by the operative hidden witness. In the exact
  recovery form used by the hidden-operator separation, \(h' = h\) for every
  hidden witness \((\Delta,h,e)\) compatible with the public observation.
\end{definition}

The exact recovery formulation is deliberately strong. It isolates the
ambiguity proved in the Lean development: the same public cochain may be
compatible with hidden operators that disagree about harmonicity. Thus a solver
that sees only public data cannot be uniformly sound unless the public
observation determines enough hidden-operator structure.

\subsection{Recovery Target}

The recovery target is the harmonic cochain selected by the operative kernel,
not necessarily the entire kernel subspace. This distinction keeps HCP parallel
to CVP: CVP asks for a closest vector in a presented lattice, while HCP asks for
a harmonic cochain in the operative Hodge kernel.

In the metric variant used for comparison with CVP, one may ask for the closest
harmonic cochain to \(b\). Once \(\Delta\) or a basis for \(\ker(\Delta)\) is
explicitly presented, this becomes a closest-vector question on the harmonic
lattice. The hidden-operator model withholds precisely that presentation.

\subsection{Computational Model}

The computational model separates three layers.
\begin{enumerate}[label=(\roman*)]
  \item The mathematical layer defines public projections, hidden operator
    fibers, harmonicity, and kernel ambiguity.
  \item External hardness assumptions and cited coding-theory theorems enter at
    explicitly labeled theorem boundaries.
  \item Runtime obligations such as byte parsing, canonical encodings, hash
    domain separation, proof verification, and state transition are outside the
    mathematical HCP definition.
\end{enumerate}

This paper focuses on the first two layers. The runtime layer is part of the
protocol and proof-carrying-verification story developed in the sequel.
